
\documentclass{article}
\usepackage{colortbl}
\usepackage{makecell}
\usepackage{multirow}
\usepackage{supertabular}

\begin{document}

\newcounter{utterance}

\centering \large Interaction Transcript for game `cladder', experiment `full\_v1.5\_default', episode 9813 with reflexion{-}Qwen3.5{-}27B.
\vspace{24pt}

{ \footnotesize  \setcounter{utterance}{1}
\setlength{\tabcolsep}{0pt}
\begin{supertabular}{c@{$\;$}|p{.15\linewidth}@{}p{.15\linewidth}p{.15\linewidth}p{.15\linewidth}p{.15\linewidth}p{.15\linewidth}}
    \# & \multicolumn{2}{c}{Player} && \multicolumn{2}{c}{Game Master} \\
    \hline

    \theutterance \stepcounter{utterance}  
    & & & \multicolumn{4}{p{0.6\linewidth}}{
        \cellcolor[rgb]{0.9,0.9,0.9}{
            \makecell[{{p{\linewidth}}}]{
                \texttt{\tiny{[P1$\langle$GM]}}
                \texttt{You have attempted this same task before and did not succeed. These are the plans you wrote after those attempts.} \\
\texttt{Trial \#1: New plan: I failed because I started with the wrong answer ("No") and then contradicted myself in the reasoning, ultimately cutting off mid{-}sentence without providing a final conclusion. The causal logic was actually correct in my draft (if \$Jyka\$ becomes true while \$Yupt\$ remains false, \$Kwox\$ must be true), but my initial commitment to "No" invalidated the entire response. Next time, I will first rigorously derive the logical outcome of the counterfactual intervention (\$do(Jyka=1)\$ given \$Yupt=0\$) to determine the correct boolean answer. I will then ensure the very first word of my response matches this derived truth value. Finally, I will provide a complete, coherent explanation that supports this answer without self{-}contradiction or truncation, ensuring the response ends with a definitive conclusion.} \\
\\ 
\texttt{Trial \#2: New plan: I failed because I prematurely committed to "No" without fully resolving the counterfactual logic, leading to a self{-}contradictory explanation that was cut off mid{-}sentence. The correct logic is that since "Jyka or Yupt causes Kwox" and we know "not Yupt," the counterfactual intervention "if Jyka instead of not Jyka" (setting Jyka to true) satisfies the condition, making Kwox true. Next time, I will first explicitly calculate the outcome of the intervention \$do(Jyka=1)\$ given \$Yupt=0\$ to confirm the answer is "Yes." I will then start my response immediately with "Yes," followed by a concise, linear derivation showing that the presence of Jyka alone is sufficient to cause Kwox in this specific scenario. I will ensure the explanation is complete and ends with a definitive conclusion, avoiding any mid{-}sentence truncation or logical backtracking.} \\
\\ 
\\ 
\texttt{You are an expert in causal inference. The following question is not a typical commonsense query, but rather a meticulously designed question created by a professor specializing in causal inference, intended to assess the students' mastery of the course content.} \\
\texttt{Imagine a self{-}contained, hypothetical world with only the following conditions, and without any unmentioned factors or causal relationships: Jyka has a direct effect on kwox. Yupt has a direct effect on kwox. We know that jyka or yupt causes kwox. We observed an individual is not yupt. Would an individual is kwox if jyka instead of not jyka?} \\
\texttt{Start your answer with 'Yes' or 'No', followed by additional reasoning or evidence to support your explanation.} \\
            }
        }
    }
     \\ \\

    \theutterance \stepcounter{utterance}  
    & \multicolumn{4}{p{0.6\linewidth}}{
        \cellcolor[rgb]{0.561,0.737,0.561}{
            \makecell[{{p{\linewidth}}}]{
                \texttt{\tiny{[P1$\rangle$GM]}}
                \texttt{Yes} \\
\\ 
\texttt{In this hypothetical causal world, the condition "Jyka or Yupt causes Kwox" establishes a logical disjunction where the presence of either \$Jyka\$ or \$Yupt\$ is sufficient to trigger \$Kwox\$. We are given the factual observation that the individual is not \$Yupt\$ (\$Yupt = 0\$). The question asks for the counterfactual outcome: what would happen to \$Kwox\$ if we intervene to set \$Jyka\$ to true (\$do(Jyka = 1)\$), replacing the actual state where \$Jyka\$ was false.} \\
\\ 
\texttt{Under the intervention \$do(Jyka = 1)\$, the condition for \$Kwox\$ becomes "True OR False" (since \$Jyka\$ is now true and \$Yupt\$ remains false). In boolean logic, if at least one of the sufficient causes is present, the effect occurs. Therefore, the intervention of making \$Jyka\$ true is sufficient to cause \$Kwox\$ to be true, regardless of the status of \$Yupt\$. Consequently, the individual would be \$Kwox\$ in this counterfactual scenario.} \\
            }
        }
    }
    & & \\ \\

    \theutterance \stepcounter{utterance}  
    & & & \multicolumn{2}{p{0.3\linewidth}}{
        \cellcolor[rgb]{0.9,0.9,0.9}{
            \makecell[{{p{\linewidth}}}]{
                \texttt{\tiny{[GM$|$GM]}}
                \texttt{Yes} \\
            }
        }
    }
    & & \\ \\

    \theutterance \stepcounter{utterance}  
    & & & \multicolumn{2}{p{0.3\linewidth}}{
        \cellcolor[rgb]{0.9,0.9,0.9}{
            \makecell[{{p{\linewidth}}}]{
                \texttt{\tiny{[GM$|$GM]}}
                \texttt{yes} \\
            }
        }
    }
    & & \\ \\

    \theutterance \stepcounter{utterance}  
    & & & \multicolumn{2}{p{0.3\linewidth}}{
        \cellcolor[rgb]{0.9,0.9,0.9}{
            \makecell[{{p{\linewidth}}}]{
                \texttt{\tiny{[GM$|$GM]}}
                \texttt{game\_result = WIN} \\
            }
        }
    }
    & & \\ \\

\end{supertabular}
}

\end{document}
